%% ============================================================
%% NEW SUBSECTION -- "Why the extended physics moves so little here"
%% Insert in Results and discussion, after the linearisation subsection and before
%% "Implications for modelling practice".
%%
%% This is the single most valuable addition identified by the revision brief, and it
%% answers R1.5 where the reviewer asked for it -- in the discussion, as a mechanism --
%% rather than only as a line in the limitations.
%%
%% THE ARGUMENT: our nulls (hydraulics <1 %, delay 0.00 %, linearisation -0.15/-0.33 %)
%% admit two readings. Either the physics genuinely does not matter, or our assumptions
%% removed the degrees of freedom through which it would act. A referee will supply the
%% second reading if we do not. We can name each closed channel and, for one of them,
%% we have already opened it and measured what happens -- which is what makes this
%% section evidence rather than hedging.
%% ============================================================

\subsection{Why the extended physics moves so little here}
\label{subsec:whynull}

Three of this paper's results are nulls: station-resolved hydraulics change dispatch cost
by under one percent, transport delay by exactly nothing, and the piecewise linearisation
by a few tenths of a percent. A null invites two readings---the phenomena are genuinely
inert, or the assumptions under which we tested them closed the channels through which they
would have acted. The second is legitimate and deserves more than a caveat.

Four assumptions each close a specific channel. A \emph{prescribed heating curve} fixes the
supply and return temperatures, so the optimiser cannot trade thermal loss against pumping
energy, the trade in which hydraulics would acquire economic weight. \emph{Precomputed
heat-pump performance} makes the coefficient of performance independent of the network
state, so a heat pump cannot respond to a locally depressed supply temperature by shifting
its output in time or place. \emph{Fixed capacities} remove siting and sizing, so spatial
structure cannot express itself through where capacity is placed, only through where heat
is routed. And \emph{hourly resolution} discretises away every transient shorter than the
time step: on this network the longest trunk travel time is about 35 minutes, so transport
delay rounds to zero steps, and the thermal inertia of the pipe volume, which buffers
precisely the sub-hourly excursions we exclude, never enters.

The honest conclusion is therefore conditional: within the deterministic, hourly,
fixed-capacity, fixed-temperature planning regime, extended network physics does not change
dispatch decisions. That regime is not a corner case---it is where most published
district-heating dispatch models operate---so the practical guidance the nulls support (spend
modelling effort on loss representation before hydraulic detail) applies wherever such models
are used.

What raises this above a plausible story is that we have opened one of the four channels
and measured the result. Freeing the supply temperature, the lever most likely to break the
hydraulic null, does break it: pipe velocity becomes the binding constraint at a 20\,K
reduction and pumping energy rises by a factor of four
(Section~\ref{subsec:tsup}). Hydraulics move from a negligible \emph{cost} to a binding
\emph{constraint} the moment the degree of freedom exists. That is direct evidence for the
mechanism rather than an appeal to it, and it sets the expectation for the three channels
we did not open: each should convert a null into a constraint rather than into a large cost
term, because the physics enters dispatch as a feasibility boundary and not as a driver of
the schedule.

One asymmetry is worth stating plainly, because it determines how far the nulls can be
trusted. Relaxing any of these four assumptions can only give the extended physics
\emph{more} room to matter, never less. The nulls are therefore conservative: they are
lower bounds on the fidelity requirement, and a reader operating outside this regime should
treat them as the floor rather than the answer. The loss-dominance result carries no such
\emph{dispatch-regime} qualification. It follows from an exact decomposition of a cost gap,
holds across 135 synthetic networks spanning two orders of magnitude in loss burden, and
does not depend on any of the four assumptions above---though, like every result here, it is
established for radial, centrally-supplied networks. That is why it, and not the nulls, is
what the title claims.
